\chapter{Conclusion}
\label{chap:conclusion}
For this study, the researchers created a computational modelling as a means to detect and quantify the existence of Braess Paradox within the Iloilo City that is based on the DPWH 2024 atlas. It provided a new methodology of evaluating the traffic management system of a complex traffic network in an urban Iloilo City.

\section{Limitations}
\label{sec:limitations}


\section{Recommendations}
\label{sec:recommendations}
Based on the findings of the study, future research should consider different networks that have accessible and available traffic data. Utilizing other networks allows comparison of accuracy of the traffic network to the reflected traffic data. High-quality data, instead of simplified ones like zonal OD matrix, would also improve the reliability of the computational model better suited for the methodology framework used by Youn et al. (2008).

Vehicles in a traffic network do not only traverse through the primary arterial roads. Further improvement of the computational modeling is to implement secondary and tertiary roads, as secondary roads will act as shortcuts. This recommendation oversees conditions where high-capacity primary roads are avoided in favor of lower-capacity secondary roads that triggers a bottleneck, therefore presenting the paradox. 

Public transport systems also occupy the traffic network while also contributing to its complexity by their unpredictable operational behavior. The inclusion of public transport in further research can help in creating a more realistic traffic network model while also considering it as a variable that may contribute to further appearance of the Braess paradox link. Since public vehicles follow fixed routes, future studies may implement them as a fixed flow that may reduce the width of the road for the private vehicles. 

Finally, discussion with stakeholders about the expansion of this study is highly recommended. The stakeholders such as the local government would help in creating a model that caters to the traffic network of Iloilo City rather than relying on standard coefficients such as the $\alpha$ and $\beta$ parameters in the BPR function. Presenting the findings to the Iloilo City Traffic and Transportation Management may help in evaluating proposed road projects.